%% CV Template: sharique
%% Recreated in LaTeX from an HTML/CSS resume export (no photo, blue accent
%% bar, two-column skills/certs grid, language grid, optional award callout).
%% See TEMPLATE.md for style rules, font notes, and known pitfalls.
%%
%% Compile with: pdflatex -interaction=nonstopmode template.tex

\documentclass[10pt,a4paper]{article}
\usepackage[top=1.1cm,bottom=1.1cm,left=1.5cm,right=1.5cm]{geometry}
\usepackage[utf8]{inputenc}
\usepackage[T1]{fontenc}
\usepackage{tgheros}
\renewcommand{\familydefault}{\sfdefault}
\usepackage[table]{xcolor}
\usepackage{graphicx}
\usepackage{tikz}
\usepackage{hyperref}
\usepackage{titlesec}
\usepackage{microtype}
\usepackage{tabularx}

\pagestyle{empty}
\setlength{\parindent}{0pt}

% ---- Colors (from the source CSS) ----
\definecolor{accent}{HTML}{185FA5}
\definecolor{graytext}{HTML}{6B7280}
\definecolor{darktext}{HTML}{1A1A1A}
\definecolor{awardbg}{HTML}{E6F1FB}
\color{darktext}

\hypersetup{colorlinks=true, urlcolor=accent, linkcolor=accent}

% Section title: bold, blue, small, thin gray rule below.
% Type section names already in UPPERCASE in the body (see usage below) -
% titlesec cannot cleanly auto-uppercase + letter-space the argument, so
% this is done manually per \section{...} call with \textls[...]{...}.
\titleformat{\section}
  {\fontsize{8}{10}\selectfont\bfseries\color{accent}}
  {}{0pt}{}
  [\vspace{2pt}{\color{black!15}\titlerule[0.4pt]}]
\titlespacing{\section}{0pt}{12pt}{8pt}

\renewcommand{\labelitemi}{\raisebox{0.15em}{\tiny$\bullet$}}

% Compact bullet list - avoids depending on enumitem (not present on every
% TeX Live install); uses plain LaTeX list-parameter overrides instead.
\newenvironment{compactitem}{%
  \begin{itemize}
  \setlength{\itemsep}{2pt}\setlength{\topsep}{2pt}%
  \setlength{\parsep}{0pt}\setlength{\partopsep}{0pt}%
  \setlength{\leftmargin}{12pt}
}{\end{itemize}}

% ---- Reusable commands ----

% Entry header row: title/name (left, bold) + date range (right, gray)
\newcommand{\entryrow}[2]{%
  \noindent\begin{tabularx}{\linewidth}{@{}X r@{}}
  {\fontsize{10}{12}\selectfont\bfseries #1} & {\fontsize{8.5}{10}\selectfont\color{graytext}#2}
  \end{tabularx}\par\vspace{1pt}}

% Entry subline: company/school (blue) - pass "\color{graytext}\textbullet\ Location"
% as part of #1 to switch to gray mid-line, matching the source design.
\newcommand{\entrysub}[1]{{\fontsize{9}{11}\selectfont\bfseries\color{accent}#1}\par\vspace{3pt}}

% Two-column grid row: bold label (line 1) + gray detail (line 2), top-aligned,
% independent minipages per row. Deliberately NOT a single shared tabularx row -
% when one cell's text wraps to more lines than its neighbor, a shared table row
% staggers every row that follows it. Two independent minipages per row-pair
% avoids that coupling entirely (used for Skills).
\newcommand{\skillrow}[4]{\noindent
  \begin{minipage}[t]{0.48\linewidth}
    {\fontsize{9}{11}\selectfont\bfseries #1}\par
    {\fontsize{8.5}{11}\selectfont\color{graytext}#2}
  \end{minipage}\hfill
  \begin{minipage}[t]{0.48\linewidth}
    {\fontsize{9}{11}\selectfont\bfseries #3}\par
    {\fontsize{8.5}{11}\selectfont\color{graytext}#4}
  \end{minipage}\par\vspace{8pt}}

% Same row-pair pattern as \skillrow, regular (non-bold) weight for the label -
% used for Certifications. Pass "\href{URL}{Name}" as a label argument for a
% linked cert (renders in accent blue via the colorlinks/urlcolor setup below -
% do NOT wrap it in \color{darktext}, that defeats the link color); pass plain
% text for an unlinked one (e.g. "in progress" - inherits darktext).
\newcommand{\certrow}[4]{\noindent
  \begin{minipage}[t]{0.48\linewidth}
    {\fontsize{9}{11}\selectfont #1}\par
    {\fontsize{8.5}{10}\selectfont\color{graytext}#2}
  \end{minipage}\hfill
  \begin{minipage}[t]{0.48\linewidth}
    {\fontsize{9}{11}\selectfont #3}\par
    {\fontsize{8.5}{10}\selectfont\color{graytext}#4}
  \end{minipage}\par\vspace{6pt}}

% Single language column: name (bold) over proficiency level (gray). Natural
% width (sized to its own content, not a fixed fraction) - matches the source
% CSS's flexbox, which does the same. Lay out with \hspace{18pt} between calls;
% see "Known pitfalls" in TEMPLATE.md if the row ever overflows \linewidth.
\newcommand{\langcol}[2]{%
  \begin{tabular}[t]{@{}l@{}}
    {\fontsize{9}{11}\selectfont\bfseries #1}\\
    {\fontsize{8.5}{10}\selectfont\color{graytext}#2}
  \end{tabular}}

% Highlighted callout box (e.g. "OKR of the Quarter") - light blue background,
% blue left bar. Built with plain colortbl (no tcolorbox dependency).
\newcommand{\awardbox}[2]{%
  \vspace{4pt}\noindent
  \begin{tabular}{@{}>{\columncolor{accent}}p{2pt}>{\columncolor{awardbg}}p{\dimexpr\linewidth-2pt-3\tabcolsep}@{}}
   & \begin{minipage}[t]{\dimexpr\linewidth-2pt-3\tabcolsep}
       \vspace{4pt}{\fontsize{9}{11}\selectfont\bfseries\color{accent}#1}\\[2pt]
       {\fontsize{8.5}{10}\selectfont #2}\vspace{4pt}
     \end{minipage}\\
  \end{tabular}\vspace{4pt}}

\begin{document}
\thispagestyle{empty}

% ============================================================
%     HEADER
% ============================================================
% No photo in the source design - full linewidth. If a specific application
% wants one, see "Adding a photo" in TEMPLATE.md's Known pitfalls.
\noindent
{\fontsize{22}{24}\selectfont\bfseries [YOUR_NAME]}\par
\vspace{2pt}
{\fontsize{11}{13}\selectfont\color{accent}[YOUR_HEADLINE]}\par
\vspace{4pt}
{\fontsize{8.5}{11}\selectfont\color{graytext}%
 [YOUR_PHONE] \quad [YOUR_EMAIL] \quad [YOUR_LOCATION] \quad
 \href{[YOUR_LINKEDIN_URL]}{\color{accent}[YOUR_LINKEDIN_DISPLAY]} \quad
 \href{[YOUR_GITHUB_URL]}{\color{accent}[YOUR_GITHUB_DISPLAY]}\\
 [YOUR_WORK_AUTH_NOTE]}

\vspace{6pt}
\noindent{\color{accent}\rule{\linewidth}{1.2pt}}
\vspace{10pt}\par

% ============================================================
%     PROFESSIONAL SUMMARY
% ============================================================
\section{\textls[150]{PROFESSIONAL SUMMARY}}
{\fontsize{9.5}{13}\selectfont [PROFILE_STATEMENT_3_TO_5_SENTENCES_WITH_BOLDED_KEY_METRICS]}

% ============================================================
%     SKILLS
% ============================================================
\section{\textls[150]{SKILLS}}
\skillrow{[SKILL_CATEGORY_1]}{[SKILL_ITEMS_1]}{[SKILL_CATEGORY_2]}{[SKILL_ITEMS_2]}
\skillrow{[SKILL_CATEGORY_3]}{[SKILL_ITEMS_3]}{[SKILL_CATEGORY_4]}{[SKILL_ITEMS_4]}
{\fontsize{9}{11}\selectfont\bfseries [SKILL_CATEGORY_FULLWIDTH]}\par
{\fontsize{8.5}{11}\selectfont\color{graytext}[SKILL_ITEMS_FULLWIDTH]}

% ============================================================
%     PROFESSIONAL EXPERIENCE
% ============================================================
\section{\textls[150]{PROFESSIONAL EXPERIENCE}}

\entryrow{[JOB_TITLE_1]}{[JOB_START]~--~[JOB_END]}
\entrysub{[COMPANY_1] \color{graytext}\textbullet\ [LOCATION_1]}
\begin{compactitem}
  \item [BULLET_WITH_\textbf{BOLDED_METRIC}]
  \item [BULLET]
  \item [BULLET]
\end{compactitem}

% Optional highlight box - use sparingly, at most one per role
\awardbox{[AWARD_TITLE]}{[AWARD_TEXT]}

\vspace{8pt}
\entryrow{[JOB_TITLE_2]}{[JOB_START]~--~[JOB_END]}
\entrysub{[COMPANY_2] \color{graytext}\textbullet\ [LOCATION_2]}
\begin{compactitem}
  \item [BULLET]
  \item [BULLET]
\end{compactitem}

% Entry with no company/location line (e.g. "Earlier Experience" roll-up):
% skip \entrysub and start bullets directly after \entryrow.

% ============================================================
%     CERTIFICATIONS
% ============================================================
\section{\textls[150]{CERTIFICATIONS}}
\certrow{\href{[CERT_URL_1]}{[CERT_NAME_1]}}{[CERT_DATE_1]}{\href{[CERT_URL_2]}{[CERT_NAME_2]}}{[CERT_DATE_2]}
\certrow{\href{[CERT_URL_3]}{[CERT_NAME_3]}}{[CERT_DATE_3]}{[CERT_NAME_4]}{[CERT_DATE_4]}

% ============================================================
%     EDUCATION
% ============================================================
\section{\textls[150]{EDUCATION}}
{\fontsize{10}{12}\selectfont\bfseries [DEGREE_1]}\par
{\fontsize{9}{11}\selectfont\color{graytext}[SCHOOL_1] \textbullet\ [DATES_1]}
\begin{compactitem}
  \item [BULLET]
\end{compactitem}

\vspace{4pt}
{\fontsize{10}{12}\selectfont\bfseries [DEGREE_2]}\par
{\fontsize{9}{11}\selectfont\color{graytext}[SCHOOL_2] \textbullet\ [DATES_2]}

% ============================================================
%     PROJECTS
% ============================================================
\section{\textls[150]{PROJECTS}}
% Wrap the title in \href when a public repo/demo URL exists - it renders in
% accent blue automatically via the colorlinks setup above. Drop the \href
% (plain title text) if there's nothing to link to.
\entryrow{\href{[PROJECT_URL]}{[PROJECT_NAME]}}{[PROJECT_START]~--~[PROJECT_END]}
\begin{compactitem}
  \item [BULLET]
  \item [BULLET]
\end{compactitem}

% ============================================================
%     OPEN SOURCE & COMMUNITY
% ============================================================
\section{\textls[150]{OPEN SOURCE \& COMMUNITY}}
\begin{compactitem}
  \item [BULLET]
  \item [BULLET]
\end{compactitem}

% ============================================================
%     LANGUAGES
% ============================================================
\section{\textls[150]{LANGUAGES}}
% Grid, not inline prose: name bold on its own line, level gray beneath, one
% natural-width column per language, spaced with \hspace.
\noindent\langcol{[LANGUAGE_1]}{[LEVEL_1]}\hspace{18pt}
\langcol{[LANGUAGE_2]}{[LEVEL_2]}\hspace{18pt}
\langcol{[LANGUAGE_3]}{[LEVEL_3]}\hspace{18pt}
\langcol{[LANGUAGE_4]}{[LEVEL_4]}

\end{document}
